\unitopen{29}
\begin{candos}
\item wyrazić życzenie, wątpliwość i nierealny warunek (\de{Wenn ich Zeit hätte, \dots});
\item dać radę (\de{Du solltest \dots}, \de{An deiner Stelle würde ich \dots});
\item utworzyć \de{Konjunktiv II} czasowników \de{sein, haben, werden, können, müssen, sollen} oraz \de{würde} + bezokolicznik;
\item mówić o przeszłości nierealnej: \de{Wenn ich das gewusst hätte, \dots}
\end{candos}

\section{Słownictwo}
\vocabtable{29}

\section{Formy}
Teraźniejszy \de{Konjunktiv II} (tryb przypuszczający) tworzysz od Präteritum. Dla większości czasowników używasz \de{würde} +
bezokolicznik; dla kilku częstych form własnych:

\begin{center}\small
\begin{tabular}{@{}llll@{}}
\toprule
\textsf{\textbf{Bezokolicznik}} & \textsf{\textbf{ich/er}} & \textsf{\textbf{du}} & \textsf{\textbf{wir/sie}}\\
\midrule
\de{sein} & \de{wäre} & \de{wärest / wärst} & \de{wären}\\
\de{haben} & \de{hätte} & \de{hättest} & \de{hätten}\\
\de{werden} & \de{würde} & \de{würdest} & \de{würden}\\
\de{können} & \de{könnte} & \de{könntest} & \de{könnten}\\
\de{müssen} & \de{müsste} & \de{müsstest} & \de{müssten}\\
\de{sollen} & \de{sollte} & \de{solltest} & \de{sollten}\\
\bottomrule
\end{tabular}
\end{center}
\textbf{Życzenie lub rada:} \de{Ich hätte gern Urlaub.} \de{Du solltest mehr schlafen.} \textbf{Nierealny warunek (teraźniejszość):}
\de{Wenn ich mehr Zeit \textbf{hätte}, \textbf{würde} ich öfter Sport machen.} \textbf{Przeszłość:} \de{hätte}/\de{wäre} + Partizip II:
\de{Wenn ich das \textbf{gewusst hätte}, \textbf{wäre} ich früher \textbf{gekommen}.}

\begin{zasada}[Dwa człony]
\de{Wenn}-zdanie opisuje warunek (czasownik na końcu), zdanie główne skutek (z \de{würde} lub formą własną): \de{Wenn ich Zeit hätte,
würde ich kommen.} Kolejność można odwrócić: \de{Ich würde kommen, wenn ich Zeit hätte.}
\end{zasada}

\begin{falle}
Polski tryb przypuszczający używa „by” i jest prosty (\emph{poszedłbym}); po niemiecku nie wolno mieszkać obu członów
z \de{würde}: \de{Wenn ich würde Zeit haben, würde ich kommen} jest błędem. Poprawnie jedna \de{würde}: \de{Wenn ich Zeit hätte,
würde ich kommen.} Drugi: używaj formy własnej dla \de{sein, haben} i modalnych: \de{Ich würde Zeit haben} brzmi sztucznie,
\de{Ich hätte Zeit}.
\end{falle}

\begin{szwajcaria}
Szwajcarzy używają \de{würde}-konstrukcji w mowie bardzo często (\de{Ich würd gern \dots}), a w dialekcie \de{wür} + infinitiv
zastępuje większość form. Łagodny ton jest normą: zamiast \de{Das ist falsch} mówi się \de{Ich wäre nicht sicher, ob das stimmt.}
\end{szwajcaria}

\begin{praca}
\de{An deiner Stelle würde ich zuerst die Logs prüfen.} \de{Wir sollten das Release verschieben.} \de{Wenn wir mehr Zeit hätten,
würden wir den Code refaktorieren.} \de{Wenn wir den Test geschrieben hätten, wäre der Fehler früher aufgefallen.}
\de{Es wäre gut, wenn wir die Deadline besprechen könnten.}
\end{praca}

\section{Dialogi}
\textbf{Dialog 1 --- rada na rozmowie}
\begin{dialog}
\turn{A}{Ich bin nervös wegen des Vorstellungsgesprächs. Was würdest du tun?}{Denerwuję się rozmową kwalifikacyjną. Co byś zrobił?}
\turn{B}{An deiner Stelle würde ich mich gut vorbereiten und die Firma anschauen.}{Na twoim miejscu dobrze bym się przygotował i obejrzał firmę.}
\turn{A}{Hätte ich mehr Zeit, würde ich länger üben. Aber ich muss auch arbeiten.}{Gdybym miał więcej czasu, ćwiczyłbym dłużej. Ale muszę też pracować.}
\turn{B}{Du solltest jeden Tag zehn Minuten sprechen. Das wäre schon viel.}{Powinieneś codziennie mówić dziesięć minut. To już dużo.}
\end{dialog}

\section{Ćwiczenia}
\begin{cwiczenia}
\exitem{Odmień Konjunktiv II: \de{ich (haben)}.}{\de{hätte}}
\exitem{Odmień Konjunktiv II: \de{er (sein)}.}{\de{wäre}}
\exitem{Odmień Konjunktiv II: \de{du (können)}.}{\de{könntest}}
\exitem{Dokończ: \de{Wenn ich Zeit \dots{} (haben), \dots{} ich kommen (würden).}}{\de{Wenn ich Zeit hätte, würde ich kommen.}}
\exitem{Dokończ: \de{Wenn ich Geld \dots{} (haben), \dots{} ich reisen (würden).}}{\de{Wenn ich Geld hätte, würde ich reisen.}}
\exitem{Daj radę koledze, który jest chory.}{\de{Du solltest zum Arzt gehen.}}
\exitem{Daj radę z formułą \de{An deiner Stelle \dots}.}{\de{An deiner Stelle würde ich zu Hause bleiben.}}
\exitem{Przeszłość nierealna: \de{Wenn ich das wissen \dots, \dots ich früher gekommen.}}{\de{Wenn ich das gewusst hätte, wäre ich früher gekommen.}}
\exitem{Popraw: \de{Wenn ich würde Zeit haben, würde ich kommen.}}{\de{Wenn ich Zeit hätte, würde ich kommen.}}
\exitem{Uprzejma prośba: \de{Können Sie das erklären?} $\rightarrow$ bardziej uprzejmie.}{\de{Könnten Sie das bitte erklären?}}
\exitem{Przetłumacz: „Chciałbym mieć więcej czasu na naukę.”}{\de{Ich wünschte, ich hätte mehr Zeit zum Lernen.}}
\exitem{Przetłumacz: „Gdybyś zapytał, pomógłbym ci.”}{\de{Wenn du gefragt hättest, hätte ich dir geholfen.}}
\end{cwiczenia}

\weekend{Napisz osiem zdań \emph{Wenn ich \dots{} hätte/wäre, würde ich \dots} o swoich marzeniach zawodowych i trzy rady dla kolegi,
który zaczyna uczyć się niemieckiego. Nagraj 90 sekund.}
